\documentclass{article}
\title{Native bee richness recovers within a decade of prairie restoration, but ground-nesting specialists lag}
\author{Miriam K. Achterberg, Rafael Quintero, Priya Natarajan, Jonah Feld}
\begin{document}\maketitle
\begin{abstract}Prairie restoration is widely promoted as pollinator habitat, yet how quickly bee communities recover is poorly known. We surveyed wild bees with pan traps and aerial netting on nine restored tallgrass prairies spanning 0-30 years since seeding and one never-plowed remnant in east-central Illinois from 2016 to 2022 (20,937 specimens, 28 taxa). Total species richness rose steeply over the first six to eight years and then levelled off, reaching values statistically indistinguishable from the remnant after roughly ten years. In contrast, ground-nesting specialist bees, including oligolectic Andrena and Calliopsis, continued to accumulate across the full chronosequence, and even the oldest restorations supported fewer specialists than the remnant. Specialist richness was positively associated with bare-ground cover at trap stations, suggesting that nesting substrate, not floral resources, limits their recovery. Management that maintains bare ground, such as frequent prescribed fire, may accelerate recovery of the species most dependent on prairie.\end{abstract}
\section{Introduction}
Less than 0.01% of Illinois's original tallgrass prairie remains. Restoration plantings on former cropland now cover far more area than remnants, and state guidance increasingly assumes these plantings will support pollinators within a few years of establishment.

Most evaluations of restored habitat measure bees in the first few years after seeding. Chronosequences, in which sites of different ages stand in for time, allow longer trajectories to be estimated, but have rarely included a remnant reference or distinguished nesting guilds.

We asked (1) how bee richness and abundance change with years since restoration; (2) whether specialist and ground-nesting bees follow the same trajectory as the community as a whole; and (3) whether floral resources or nesting substrate better explain the recovery of specialists. \cite{winfree2009,kremen2002}
\section{Methods}
Study sites. Ten sites in Champaign, Piatt, Vermilion, Douglas, Macon, Coles and Ford counties were sampled. Nine are restorations seeded between 1992 and 2019 with local-ecotype mixes; one is a never-plowed remnant on private land whose location is withheld here at the request of the Illinois Department of Natural Resources because it supports the federally endangered rusty patched bumble bee.

Bee sampling. At three stations per site we set fifteen 3.25 oz bowls (five each of fluorescent blue, fluorescent yellow and white) filled with soapy water for 24 hours, four times per season (May-August). Bowls found tipped or flooded were recorded as zero effort. Netting along a 50 m transect for 20 minutes per station supplemented pan traps.

Identification. Specimens were pinned and identified to species where possible using regional keys; 15% were DNA-barcoded (COI) to resolve Lasioglossum and Melissodes. Vouchers are deposited in the university insect collection.

Analysis. Rarefied richness was modelled with generalised additive mixed models with a smooth of restoration age, site as a random intercept, and log bowl-days as an offset. Specialist richness was modelled separately, with and without the remnant site, and with bare-ground cover and floral density as covariates.
\bibliographystyle{apalike}
\bibliography{refs}
\end{document}
